\begin{proof}
We use the curve constructed in \cite[\S17]{GJR26}.  The ground field $K$ is
complete, algebraically closed, and nontrivially valued, with algebraically
closed residue field not algebraic over a finite field.  The metrized complex
in that construction has a theta graph with a vertex $P$.  The residual
curve at $P$ is an elliptic curve $E$; the three graph directions $a,b,c$ at
$P$ correspond to distinct points $x,y,z\in E(\widetilde K)$, where $z$ is
the identity and $x$ is non-torsion.  The realization theorem used there
gives a smooth proper connected genus-$3$ curve $C/K$ with this metrized
complex.  Write $\tau:C^{\mathrm{an}}\to\Gamma$ for the retraction onto its
theta skeleton.

The proof of \cite[Proposition~17.4]{GJR26} subdivides the theta skeleton to
$\Gamma'$ and constructs a piecewise $\mathbf Z$-linear function
$F:\Gamma'\to\mathbf R$, affine on every edge, whose outgoing slopes at $P$
in the directions $a,b,c$ are respectively $1,-1,0$.  It also constructs a
vertical Cartier divisor $D$ on the corresponding strictly semistable model
such that
\[
                         H:=F\circ\tau=-\log\lVert1\rVert_{\mathcal O(D)}
\]
on $C^{\mathrm{an}}$.  Thus $H$ is a globally defined integral PL function;
we do not assert that it is harmonic away from the star chosen below.  The
same proof computes the restriction of $\mathcal O(D)$ to the residual
elliptic curve as the class of $[z]-[x]$.

Choose type~$2$ points $Q_a,Q_b,Q_c$ at value-group-rational positions in the
interiors of the three graph edges and put
\[
              \Sigma_0=\bigcup_{e\in\{a,b,c\}}[P,Q_e],
              \qquad Y_0=\tau^{-1}(\Sigma_0).
\]
The restriction $F|_{\Sigma_0}$ is affine on each arm, and its outgoing
slopes $1,-1,0$ sum to zero at the sole interior vertex $P$; no balancing
condition is imposed at the boundary points $Q_e$.  Therefore
\cite[Proposition~15.9]{GJR26} makes $H|_{Y_0}$ integral harmonic.  The
domain $Y_0$ is compact and strictly affinoid, with analytic boundary
$\{Q_a,Q_b,Q_c\}$.  Put
\begin{equation}
\label{eq:ch19-harmonic-reg-open-star}
                         U=Y_0\setminus\partial Y_0,
\end{equation}
and set $h=H|_U$.  Then $U$ is a relatively compact connected open
neighbourhood of $P$, and $h$ is integral PL and harmonic.

Its residual slope divisor $D_P(h)$ is supported on
\begin{equation}
\label{eq:ch19-harmonic-reg-support}
                         S=\{x,y,z\}\subset E(\widetilde K),
\end{equation}
and its class is
\begin{equation}
\label{eq:ch19-harmonic-reg-nonzero-class}
 [D_P(h)]=[z]-[x]\neq0
       \quad\text{in}\quad
 \operatorname{Pic}^0(E)\otimes_{\mathbf Z}\mathbf R.
\end{equation}
Indeed, $[z]-[x]$ has infinite order by the choice of $x$.  Its infinite
cyclic subgroup remains injective after tensoring with $\mathbf R$, because
$\mathbf R$ is torsion-free, hence flat, over $\mathbf Z$.  By
\cite[Proposition~10.17]{GJR26}, $h$ is weakly smooth and
$\dprimeddprime h=0$.

Choose type~$2$ points $R_e\in(P,Q_e)$ at value-group-rational positions, put
\[
              \Sigma_Y=\bigcup_{e\in\{a,b,c\}}[P,R_e],
              \qquad Y=\tau^{-1}(\Sigma_Y)\Subset U,
\]
and subdivide the stable model at $R_a,R_b,R_c$.  The corresponding formal
open gives a strictly semistable formal model $\mathfrak Y$ of the strict
affinoid domain $Y$, whose skeleton is $\Sigma_Y$ and whose boundary is
$\{R_a,R_b,R_c\}$.  Let $g$ be harmonic on $U$ in the sense of
Gubler--Jell--Rabinoff.  By \cite[Proposition~15.7]{GJR26}, $g|_Y$ is
harmonic in Thuillier's sense.  Thuillier's strict-affinoid skeletal
description \cite[Definitions~2.3.1, 2.3.12, and 2.3.16]{Thu05} places
$g|_Y$ in $H(Y)$.  The independence of this space from the chosen strictly
semistable model \cite[Proposition~2.3.3]{Thu05} gives
$g|_Y=G\circ\tau_{\mathfrak Y}$, where $G$ is piecewise affine on
$S(\mathfrak Y)$ and has zero sum of outgoing slopes away from $\partial Y$.
At every subdivision point in the interior of an arm, the two one-sided
slopes therefore agree.  Thus $G$ is affine with constant slope on each
$[P,R_e]$.  Since $g$ factors through this three-armed retraction, all other
tangent directions at $P$ have slope zero.  Consequently
\begin{equation}
\label{eq:ch19-harmonic-reg-fixed-support}
       D_P(g)\in
       V_S:=\operatorname{Div}^0_S(E)\otimes_{\mathbf Z}\mathbf R,
\end{equation}
and the three coefficients of $D_P(g)$ are the outgoing slopes of $g$ on the
fixed arms of $Y$.  Here $D_P(g)$ is the residual slope divisor; its class is
the residue class of the germ of $g$ at $P$.

Suppose in addition that $g$ is ordinary CLD-smooth.  By
\cite[Remark~15.10]{GJR26}, its germ lies in the ordinary smooth harmonic
subsheaf $\mathcal F_P$.  The description of
$(\mathcal H/\mathcal F)_P$ by \cite[Lemma~2.3.22]{Thu05} then says that its
residual divisor is a real linear combination of principal divisors, or
equivalently that its class in
$\operatorname{Pic}^0(E)\otimes_{\mathbf Z}\mathbf R$ vanishes.  Together
with \eqref{eq:ch19-harmonic-reg-fixed-support}, this gives
\begin{equation}
\label{eq:ch19-harmonic-reg-picard-kernel}
 D_P(g)\in
 P_S:=\ker\!\left(
 V_S\longrightarrow\operatorname{Pic}^0(E)\otimes_{\mathbf Z}\mathbf R
 \right).
\end{equation}
The space $V_S$ is finite-dimensional and $P_S$ is a linear subspace, hence
$P_S$ is closed in the coefficient topology of $V_S$.

Suppose that a net $(h_i)$ of ordinary CLD-smooth harmonic functions on $U$
converges locally uniformly to $h$.  The inclusion of ordinary smooth forms
in weakly smooth forms, together with \cite[Propositions~10.14 and
10.17]{GJR26}, shows that every $h_i$ is harmonic in the
Gubler--Jell--Rabinoff sense.  The fixed-support argument makes $h_i$ and $h$
affine on each $[P,R_e]$.  If $\ell_e$ is the length of this segment, the
outgoing slope is
\begin{equation}
\label{eq:ch19-harmonic-reg-slope-values}
                         \frac{h_i(R_e)-h_i(P)}{\ell_e},
\end{equation}
and the analogous formula holds for $h$.  Uniform convergence on $Y$
therefore gives coefficientwise convergence
\begin{equation}
\label{eq:ch19-harmonic-reg-divisor-limit}
                         D_P(h_i)\longrightarrow D_P(h)
                         \quad\text{in }V_S.
\end{equation}
Every term on the left belongs to the closed subspace $P_S$ by
\eqref{eq:ch19-harmonic-reg-picard-kernel}.  Hence $D_P(h)\in P_S$, contrary
to \eqref{eq:ch19-harmonic-reg-nonzero-class}.

Finally, suppose that $(h_i)$ is a decreasing net of ordinary CLD-smooth
harmonic functions on $U$ with pointwise limit $h$.  On the compact set $Y$,
the net form of Dini's theorem makes the convergence uniform.  Indeed, for
$\varepsilon>0$ the increasing family of open sets
\[
              \{p\in Y:h_i(p)-h(p)<\varepsilon\}
\]
covers $Y$; compactness supplies finitely many indices and directedness one
index dominating them all.  The preceding contradiction applies.
\end{proof}
